\documentclass{article}%
\usepackage[T1]{fontenc}%
\usepackage[utf8]{inputenc}%
\usepackage{lmodern}%
\usepackage{textcomp}%
\usepackage{lastpage}%
\usepackage{geometry}%
\geometry{margin=2.2cm}%
%
\usepackage{amsmath}%
\usepackage{amssymb}%
\usepackage{booktabs}%
\usepackage{longtable}%
\usepackage{array}%
\usepackage{hyperref}%
\hypersetup{hidelinks}%
\title{Soluções\\[4pt]\large Comparação com o gabarito Gabarito preliminar oficial (SSPM, 01/09/2026), prova AMARELA\thanks{Gerado com gpt-resolve: \url{https://github.com/lgabs/gpt-resolve}}}%
\date{29/10/2026}%
%
\begin{document}%
\normalsize%
\maketitle%
\section*{Resumo}\textbf{Nota média: 10,0/10} (22 de 22 itens com a resolução correta; cada item vale o mesmo dentro da questão).\par\medskip%
\begin{longtable}{@{}c c p{0.78\textwidth}@{}}\toprule Questão & Nota & Resumo \\ \midrule
19 & 10,0 & A resolução e o gabarito preliminar oficial batem: os dois marcam (D), 1/(4$\alpha$1), e essa resposta está correta. Para os centros coincidirem, os raios finais precisam ser iguais. Como $\alpha$2 = 0,8$\alpha$1, a temperatura final se cancela e sobra 0,8$\alpha$1$\Delta$T = 0,2. Conferi a conta à mão e numericamente, e o enunciado transcrito está fiel ao PDF. Só ficaram observações menores: o resultado vale porque o modelo de dilatação é linear, o $\Delta$T é irrealista para metais e um comentário qualitativo sobre o sinal está impreciso. \\
20 & 10,0 & A resolução e o gabarito oficial preliminar (SSPM, prova AMARELA) coincidem: alternativa (E), t = $\surd$(2HM/F). Refiz a física e ela está correta. A aceleração relativa entre o piso e o corpo é (F/M - g) - (-g) = F/M, e o corpo parte do repouso em relação à cabine, então H = $\frac12$(F/M)t$^2$. O resultado não depende de m nem da velocidade inicial comum. Não há divergência a levar ao professor, só hipóteses implícitas razoáveis: F age só sobre a cabine e continua constante depois do rompimento. \\
21 & 10,0 & A resolução e o gabarito oficial preliminar marcam a mesma alternativa, (A), e ela está correta. Pela figura original, a espira externa é percorrida no sentido horário e a interna no anti{-}horário, então os campos delas em O têm sentidos opostos. O fio retilíneo contribui com $\mu$0I/(2$\pi$R), no mesmo sentido do campo da espira externa. O campo resultante se anula em um único raio, r* = $\pi$R/($\pi$+1) $\approx$ 0,76R, o que torna I verdadeira. A espira interna domina quando r$\to$0, então II é verdadeira. III é falsa porque o campo do fio não depende de r, e IV é falsa porque as espiras têm sentidos opostos. Não há divergência a resolver. \\
22 & 10,0 & A resolução e o gabarito oficial preliminar indicam a mesma alternativa, (A), e as duas estão certas. Em t = 3$\pi$ s, o disco A girou 3$\pi$/2 rad no sentido anti{-}horário: P está no topo e a seta de A aponta para baixo. O disco B girou 3$\pi$ rad no sentido horário, o que equivale a meia volta, e sua seta aponta para a direita. A física e as contas da resolução estão corretas. O único problema é a transcrição das alternativas (C) e (E): em (E), a seta de B aponta para a esquerda, não para baixo. Por isso o argumento usado para descartar (E) está impreciso, mas a conclusão vale do mesmo jeito. \\
23 & 10,0 & A resolução e o gabarito batem: os dois dão a alternativa B, 0,50 V. No gráfico, R0 = 10 $\Omega$ e Pmax = 0,025 W. Daí U = $\surd$(Pmax$\cdot$R) = $\surd$0,25 = 0,50 V, que é metade da fem ($\varepsilon$ = 1,0 V). Conferi a conta de forma independente, inclusive com o ponto em R = 50 $\Omega$, e o resultado é coerente com a curva. Não encontrei erro em nenhum dos dois. \\
24 & 10,0 & A resolução e o gabarito oficial preliminar (prova AMARELA) dão a mesma resposta: (E). Refiz o equilíbrio da barra (momentos em A e resultante nula) e obtive F\_x = 4Mg sen$\theta$ cos$\theta$ > 0 e F\_y = Mg(1 - 4cos$^2$$\theta$) < 0, já que $\theta$ < 60$^\circ$ na figura. A força da articulação aponta para a direita e para baixo, ao longo da linha 1. A resolução está correta, bem justificada e confere o resultado pelo teorema das três forças. Não há erro em nenhum dos dois. As dúvidas listadas são só observações menores sobre hipóteses e a leitura da figura. \\
25 & 10,0 & A resolução e o gabarito preliminar oficial (D) concordam, e ambos estão corretos. A reta mais inclinada no gráfico Q$\times$T corresponde a Cp, então A é o processo isobárico, e nele W = nR(T2 - T1) > 0 independentemente da atomicidade do gás. As outras alternativas foram descartadas corretamente: p2 < p3, existe trabalho no processo A, $\Delta$U depende de Cv e U2 = U3. Não encontrei erro de física nem de conta. Ficam só observações menores de redação ou hipóteses implícitas. \\
26 & 10,0 & A resolução e o gabarito preliminar oficial (prova AMARELA) concordam: alternativa C, razão 0,5. Refiz as contas e elas batem: $\eta$\_M = 1 - 1/(5 ln(5/4)) $\approx$ 0,10 e $\eta$\_máx = 0,20, o que dá razão $\approx$ 0,5 (sem arredondar, 0,519). A física da resolução está correta. Ela usa a 1ª lei em cada etapa, confere a expressão da área pela adiabática e faz o balanço de trabalho do ciclo. Não encontrei erros. Só ficam como implícitas as hipóteses de gás ideal e de diatômico sem vibração. \\
27 & 10,0 & A resolução está correta e bate com o gabarito oficial preliminar: alternativa (A). Não há força externa resultante no plano horizontal e a mola é interna, então o CM faz movimento retilíneo uniforme. Ele parte de (2; 1,5) m com velocidade (0,5; 1,5) m/s, o que dá y = 3x - 9/2. O descarte das outras alternativas está bem justificado. Só há um detalhe de digitação sem importância. \\
28 & 10,0 & A resolução e o gabarito batem e estão certos: a resposta da questão 28 é a (B), E = -($\pi$$^2$Mk$^2$q$^{4}$/T$^2$)\^{}\{1/3\}. Refiz a conta com o raio d/2 de cada partícula em torno do CM e com E = -kq$^2$/d (virial), e conferi numericamente. O único problema é um comentário secundário da resolução: ela diz que a (D) sai de usar o raio d no lugar de d/2, mas esse erro sozinho dá fator 2 e não 4. Isso não afeta a resposta. \\
29 & 10,0 & Questão 29: a resolução e o gabarito oficial preliminar (prova AMARELA) dão a mesma resposta, (A) F = GM$^2$/(10R$^2$), e ela está correta. A superposição dá F = (1/4 - 1/50)$\cdot$GMm/R$^2$ = (23/100)$\cdot$GMm/R$^2$. Com m = 10M/23, o fator 23 se cancela e resta 1/10. Conferi o enunciado no PDF original e a conta com python. \\
30 & 10,0 & A resolução escolhe (D) 6,0$\times$10\^{}{-}7 m e o gabarito preliminar oficial também é D, então resolução e gabarito estão corretos. A física está certa: franjas claras em y\_n = n$\lambda$L/d, D é a 3ª franja e vale y3 = 3$\lambda$L/d. A leitura do gráfico também está certa: $\Delta$L = 500d para $\Delta$y3 = 9,0$\times$10\^{}{-}4 m, o que dá $\lambda$ = 600 nm. Há só observações menores sobre o comentário a respeito dos distratores. \\
31 & 10,0 & A resolução e o gabarito oficial preliminar (prova AMARELA) dão a mesma resposta, (C) v = 0,20 m/s, e ambos estão certos. A resolução leva em conta as duas causas da variação do fluxo: o campo B que cresce com o tempo e a área da espira que aumenta. Assim, $\varepsilon$ = 0,040 + 0,40vt, e em t = 2,0 s a condição $\varepsilon$ = RI = 0,20 V dá v = 0,20 m/s. Refiz a conta, ela confere numericamente, e a transcrição bate com o PDF original. Não há erro na resolução nem no gabarito. \\
32 & 10,0 & A resolução e o gabarito oficial preliminar (SSPM, prova AMARELA) dão a mesma resposta, (E), e ela está correta. Pela figura original, o fundamental é a curva A, com período 2t' e f0 = 1/(2t'). B é o 3º harmônico e C é o 5º, então f\_B/f\_C = 3/5. Daí vêm v = 2L$\cdot$f0 = 0,60/t' e T/$\mu$ = (0,60/t')$^2$. As outras alternativas são falsas. Não há divergência nem erro; as únicas ressalvas, de baixa importância, são a leitura de um gráfico esquemático e a fase das componentes, que a resolução assume sem dizer. \\
33 & 10,0 & A resolução e o gabarito preliminar oficial marcam a mesma alternativa (C), e os dois estão certos. Refiz a superposição das ondas: o som percebido tem frequência (f1+f2)/2 = f2 + 20 Hz, os batimentos têm frequência |f1-f2| = 40 Hz e aparecem em qualquer posição de O, porque a posição só muda uma fase constante. As alternativas A, B, D e E foram descartadas com justificativa física correta. Não encontrei erro nem divergência. \\
34 & 10,0 & A resolução está certa e o gabarito preliminar, não. Pelo gráfico, k = 20 N/m sem ambiguidade (a curva passa em (0,10 m; 0,10 J)). Como o atrito só acelera m2, a deformação máxima sem deslizamento é d\_máx = $\mu$e$\cdot$m1$\cdot$g$\cdot$(m1+m2)/(k$\cdot$m2) = 30 cm, e não 20 cm nem 80 cm. O período é T = 2$\pi$$\surd$(3/20) $\approx$ 2,43 s, diferente de 0,4$\pi$$\surd$5 $\approx$ 2,81 s. Com os dados impressos, as três afirmativas são falsas e a resposta é B. O gabarito E parece combinar o erro clássico k$\cdot$d = $\mu$e$\cdot$m1$\cdot$g com um período que não sai dos dados. Curiosamente, com m1 = 2 kg a física correta daria exatamente a alternativa A, o que sugere um possível problema na elaboração da questão. Vale verificar se o gabarito definitivo alterou ou anulou a questão. \\
35 & 10,0 & A resolução e o gabarito oficial preliminar concordam: a resposta é a alternativa D, h\_min = 2,0 m. A física está correta. Pelo Bernoulli com velocidade uniforme no tubo, a pressão interna é mínima no ponto mais alto B, onde p0 - p\_B = $\rho$g(H+a-h). A condição $\leq$ 25 kPa dá h $\geq$ 4,5 - 2,5 = 2,0 m. Conferi a transcrição com o original e refiz as contas numericamente. Não há erro relevante, apenas uma hipótese implícita razoável (tubo de seção uniforme). \\
36 & 10,0 & A resolução está certa e concorda com o gabarito oficial preliminar: a resposta é (B). Conferi o enunciado com o PDF original. Com a placa metálica neutra entre as placas, o sistema equivale a dois capacitores em série com a mesma carga. O campo é igual nas duas regiões de ar, E = V/(d1+d2), e portanto V1/V2 = d1/d2, com V1 + V2 = V. Não há erro de física nem de conta, e não há ponto que possa mudar a nota. \\
37 & 10,0 & A resolução está correta e coincide com o gabarito preliminar oficial (B). A componente x da normal, mg$\cdot$cos30$^\circ$$\cdot$sen30$^\circ$, tem o mesmo sentido do deslocamento horizontal d$\cdot$cos30$^\circ$. Daí W = mgh$\cdot$cos$^2$30$^\circ$ = 0,60 J. Esse valor se cancela com o trabalho de N\_y (-0,60 J), e o trabalho total da normal fica nulo. Não há erro nem ambiguidade que mude a resposta. \\
38 & 10,0 & A resolução e o gabarito batem e ambos estão certos: a resposta é (C), x = (15/8)$\cdot$mg/k. Refiz a conta: o empuxo sobre a metade submersa vale mg/16, a tração no fio vale 15mg/16 e, pelos torques em torno de O, kx = 2T = 15mg/8. O enunciado confere com o PDF original. O único ponto a observar é uma hipótese já declarada na resolução e confirmada pela figura: o cilindro não toca o fundo. \\
39 & 10,0 & A resolução e o gabarito oficial preliminar (SSPM, prova AMARELA) marcam a mesma alternativa, (E), e os dois estão certos. A física da resolução está correta: com fluxo de calor constante, o calor cedido por A no intervalo todo vale P$\cdot$t2, e o calor sensível recebido por B depois da fusão vale P(t2 - t1). Dividindo um pelo outro, sai a razão da alternativa (E). As verificações de unidade e do caso t1 = 0 também estão certas, assim como a eliminação das outras alternativas. Fica só uma observação menor sobre o enunciado: C\_B deveria ser a capacidade térmica de B líquida. \\
40 & 10,0 & Questão 40: a resolução e o gabarito oficial preliminar marcam a mesma alternativa, (D) $\surd$6/2, e ela está correta. Como o barco anda só para leste, a componente norte do vento é igual nos dois referenciais, e daí v\_vb/v\_vm = cos30$^\circ$/cos45$^\circ$ = $\surd$6/2. A componente leste confirma uma solução física com v\_vm = 10($\surd$3+1) $\approx$ 27,3 km/h e v\_vb $\approx$ 33,5 km/h; conferi esses valores numericamente. A figura e o texto do original batem com a transcrição. \\
\bottomrule\end{longtable}%
\section*{Principais pontos de dúvida}%
\begin{itemize}%
\item%
\textbf{Q34} \textit{(key)} O gabarito preliminar (E) marca como verdadeiros d\_máx = 20 cm e T = 0,4$\pi$$\surd$5. Com os dados impressos (k = 20 N/m pelo gráfico, m1 = 1 kg, m2 = 2 kg, $\mu$e = 0,4), d\_máx = 30 cm e T = 2$\pi$$\surd$0,15 $\approx$ 2,43 s $\neq$ 2,81 s. O valor de 20 cm sai da análise errada k$\cdot$d = $\mu$e$\cdot$m1$\cdot$g, que ignora que o atrito só precisa acelerar m2. O período 0,4$\pi$$\surd$5 exigiria (m1+m2)/k = 0,2. Tudo indica erro no gabarito preliminar, e vale conferir se o definitivo alterou ou anulou a questão.%
\item%
\textbf{Q34} \textit{(statement)} Possível erro de digitação no enunciado: com m1 = m2 = 2,0 kg (M = 4 kg) e k = 20 N/m, a física correta dá d\_máx = 0,4$\cdot$20$\cdot$4/(20$\cdot$2) = 0,80 m = 80 cm e T = 2$\pi$$\surd$(4/20) = 0,4$\pi$$\surd$5, ou seja, a alternativa A (V)(F)(V), exatamente. Isso sugere que a questão pode ter sido elaborada com outros valores de massa. Mesmo assim, nenhuma combinação natural de massas leva à E com a análise correta. Com os dados impressos, a resposta é B.%
\end{itemize}%
\section*{Questão 19}%
\begin{longtable}{@{}p{0.06\textwidth} p{0.33\textwidth} p{0.33\textwidth} p{0.1\textwidth} p{0.08\textwidth}@{}}\toprule Item & Resolução & Gabarito & Veredito & Certo \\ \midrule
19 & (D) T\_\{i,2\} - T\_\{i,1\} = 1/(4$\cdot$$\alpha$1) & D & igual & ambos \\
\bottomrule\end{longtable}%
\begin{itemize}%
\item%
\textbf{19)} Conferi o enunciado no PDF (prova AMARELA, p. 8) e a transcrição está fiel: 0,8R para o anel 1, R para o anel 2 e $\alpha$2/$\alpha$1 = 0,8. Para os centros coincidirem, os raios finais precisam ser iguais: 0,8R{[}1+$\alpha$1(Tf-Ti1){]} = R{[}1+0,8$\alpha$1(Tf-Ti2){]}. Como 0,8R$\cdot$$\alpha$1 = R$\cdot$$\alpha$2, o termo em Tf some dos dois lados e sobra 0,8$\alpha$1(Ti2-Ti1) = 0,2, ou seja, $\Delta$T = 1/(4$\alpha$1). Testei numericamente com vários valores de Tf e a diferença entre os raios dá zero em todos. As outras opções dão 0,8$\alpha$1$\Delta$T igual a 3,2 (A), 0,8 (B), 3,6 (C) e 0,4 (E), nenhuma igual a 0,2. A resolução e o gabarito estão certos.%
\end{itemize}%
\paragraph{Pontos de dúvida}%
\begin{itemize}%
\item%
\textit{(statement)} O valor de Tf só desaparece da conta porque o modelo é linear e cada raio é medido na temperatura inicial do próprio anel. Se o cálculo usasse uma temperatura de referência comum ou a dilatação exponencial, Tf não se cancelaria exatamente. A resolução usa o modelo de livro{-}texto, que é claramente o pretendido, mas não comenta essa dependência.%
\item%
\textit{(statement)} O enunciado não é realista: com $\alpha$ \textasciitilde{} 1e{-}5 K$^{-}$$^{1}$, a resposta dá $\Delta$T \textasciitilde{} 2,5$\cdot$10$^{4}$ K, e a variação de 20\% no raio fica muito fora do regime em que $\alpha$$\Delta$T ? 1. A própria resolução aponta isso. Não muda a alternativa.%
\item%
\textit{(solution)} Na verificação de sinal, a resolução diz que o anel 1 ganha raio em relação ao anel 2 'quando ele se aquece enquanto o anel 2 esfria'. Isso é impreciso. A diferença entre os raios não depende de Tf, e o ganho relativo existe até se os dois aquecerem ou os dois esfriarem, porque o anel 1 sempre passa por uma variação de temperatura maior em Ti2 - Ti1. A observação é só qualitativa e não afeta a resposta.%
\end{itemize}%
\section*{Questão 20}%
\begin{longtable}{@{}p{0.06\textwidth} p{0.33\textwidth} p{0.33\textwidth} p{0.1\textwidth} p{0.08\textwidth}@{}}\toprule Item & Resolução & Gabarito & Veredito & Certo \\ \midrule
20 & (E) t = $\surd$(2HM/F) & E & igual & ambos \\
\bottomrule\end{longtable}%
\begin{itemize}%
\item%
\textbf{20)} Conferi o enunciado no PDF original (pág. 8/18, prova AMARELA), e a transcrição está fiel. Refiz a conta. Quando o fio se rompe, as velocidades do corpo e da cabine são iguais. Depois disso, o corpo cai livremente (a = -g) e a cabine passa a ter aceleração a = (F - Mg)/M = F/M - g. A aceleração relativa é F/M, com velocidade relativa inicial nula. Então H = $\frac12$(F/M)t$^2$, o que dá t = $\surd$(2HM/F), alternativa (E). As verificações da resolução estão corretas: as unidades dão s$^2$, e no caso F = Mg o resultado se reduz a $\surd$(2H/g). A análise das alternativas erradas também está correta. A resolução e o gabarito coincidem.%
\end{itemize}%
\paragraph{Pontos de dúvida}%
\begin{itemize}%
\item%
\textit{(statement)} O enunciado não diz que F atua apenas sobre a cabine, nem que continua igual depois do rompimento. A resolução declara essas duas hipóteses, e elas são a leitura natural ("força externa constante"). Se F mudasse com o rompimento, o resultado seria outro, mas nada no texto sugere isso.%
\item%
\textit{(solution)} A resolução diz que (C) e (D) "somam ou subtraem g e F/M de forma incorreta". A justificativa é um pouco vaga, mas não afeta a resposta: (C) é dimensionalmente consistente e só erra a aceleração relativa.%
\end{itemize}%
\section*{Questão 21}%
\begin{longtable}{@{}p{0.06\textwidth} p{0.33\textwidth} p{0.33\textwidth} p{0.1\textwidth} p{0.08\textwidth}@{}}\toprule Item & Resolução & Gabarito & Veredito & Certo \\ \midrule
21 & (A) Apenas I e II são verdadeiras. Em O, B = $\mu$0I/(2R)(1+1/$\pi$) - $\mu$0I/(2r), que se anula só em r* = $\pi$R/($\pi$+1) $\approx$ 0,76R; III e IV são falsas. & A & igual & ambos \\
\bottomrule\end{longtable}%
\begin{itemize}%
\item%
\textbf{21)} Refiz a questão a partir da figura ampliada da prova. O fio retilíneo da esquerda entra no arco superior direito da espira externa e percorre essa espira no sentido horário: as setas apontam para a direita no topo esquerdo, para baixo/direita no topo direito e para a esquerda embaixo. Em seguida o fio desce pelo trecho radial esquerdo e percorre a espira interna no sentido anti{-}horário (seta embaixo apontando para a direita). Depois sobe pelo trecho radial direito (seta para cima) e sai pelo fio retilíneo da direita.\newline%
\newline%
Contribuições em O:\newline%
{-} espira externa: $\mu$0I/(2R), entrando;\newline%
{-} fio retilíneo (corrente para a direita, O abaixo, distância R): $\mu$0I/(2$\pi$R), entrando;\newline%
{-} espira interna: $\mu$0I/(2r), saindo.\newline%
\newline%
Portanto B(r) = $\mu$0I/(2R)(1+1/$\pi$) - $\mu$0I/(2r). Essa função é estritamente crescente em r e só se anula em r* = $\pi$R/($\pi$+1) $\approx$ 0,7585R, que está dentro de (0, R). Logo I é verdadeira. Quando r$\to$0, o termo 1/r domina, então II é verdadeira. O campo do fio retilíneo só depende de I e R, então III é falsa. As espiras têm sentidos opostos, então IV é falsa. A resposta é (A), igual à da resolução e à do gabarito.%
\end{itemize}%
\paragraph{Pontos de dúvida}%
\begin{itemize}%
\item%
\textit{(statement)} A transcrição diz que a seta da espira interna fica 'no lado direito'. No original ampliado, a seta de corrente da espira interna está na parte de baixo e aponta para a direita (sentido anti{-}horário). A seta no lado inferior direito é a ponta do raio tracejado r. A resolução leu a figura corretamente.%
\item%
\textit{(solution)} A resolução trata cada espira como uma circunferência completa e o trecho retilíneo como dois fios semi{-}infinitos terminando no ponto de tangência, ou seja, um fio infinito à distância R. Para isso, despreza as pequenas aberturas dos trechos radiais e a pequena curva de junção no topo. Essa é a idealização que o enunciado pretende.%
\item%
\textit{(key)} A afirmativa I só é verdadeira se a contribuição do fio retilíneo for incluída. Contando só as duas espiras, B = $\mu$0I/2$\cdot$(1/R - 1/r) se anularia apenas em r = R, valor fora do intervalo aberto 0 < r < R. O gabarito (A) é coerente com a inclusão do fio, como o enunciado pede ao mencioná{-}lo na afirmativa III.%
\end{itemize}%
\section*{Questão 22}%
\begin{longtable}{@{}p{0.06\textwidth} p{0.33\textwidth} p{0.33\textwidth} p{0.1\textwidth} p{0.08\textwidth}@{}}\toprule Item & Resolução & Gabarito & Veredito & Certo \\ \midrule
22 & (A): P no topo de A, seta de A para baixo, seta de B para a direita ($\Delta$$\theta$\_A = 3$\pi$/2 rad no sentido anti{-}horário; $\Delta$$\theta$\_B = 3$\pi$ rad no sentido horário) & A & igual & ambos \\
\bottomrule\end{longtable}%
\begin{itemize}%
\item%
\textbf{22)} Refiz as contas (inclusive em python3). Inclinação: tan 30$^\circ$ = $\surd$3/3 m/s$^2$. V = $\surd$3$\cdot$$\pi$ m/s. Arco percorrido = área sob o gráfico = 3$\surd$3$\pi$$^2$/2 m. Logo $\Delta$$\theta$\_A = 3$\pi$/2 (270$^\circ$ no sentido anti{-}horário): a seta de A vai de 0$^\circ$ para 270$^\circ$, ou seja, aponta para baixo, e P vai de 180$^\circ$ para 90$^\circ$, ficando no topo. Sem deslizamento, $\Delta$$\theta$\_B = (R\_A/R\_B)$\cdot$$\Delta$$\theta$\_A = 3$\pi$, no sentido horário. Isso equivale a meia volta, e a seta de B passa de 180$^\circ$ para 0$^\circ$, apontando para a direita. No original, a alternativa (A) mostra exatamente essa configuração. A resolução e o gabarito concordam e estão corretos.%
\end{itemize}%
\paragraph{Pontos de dúvida}%
\begin{itemize}%
\item%
\textit{(statement)} A transcrição das alternativas (C) e (E) não confere com o PDF original. Em (E), a seta de B aponta para a ESQUERDA, para o ponto de contato, e não para baixo. Em (C), a seta de B é visível e aponta para a direita. A resolução descartou (E) com base na leitura errada ('seta de B para baixo $\Rightarrow$ $\Delta$$\theta$\_B = $\pi$/2'). No original, (E) corresponde a $\Delta$$\theta$\_B $\equiv$ 0 (mod 2$\pi$) e continua incompatível com $\Delta$$\theta$\_B = 3$\pi$. A justificativa desse descarte está imprecisa, mas a resposta não muda.%
\item%
\textit{(solution)} A resolução lê a inclinação como tan 30$^\circ$ em (m/s)/s. Isso só vale se os dois eixos estiverem na mesma escala numérica, o que é a convenção implícita da questão. No desenho, a reta parece mais inclinada que 30$^\circ$, mas o rótulo prevalece. Com o ângulo de 30$^\circ$ medido a partir do eixo v (inclinação $\surd$3), P terminaria embaixo, configuração que não aparece em nenhuma alternativa. Isso confirma a leitura usada.%
\item%
\textit{(statement)} O enunciado dá R\_A = $\surd$3$\pi$ sem unidade; a resolução supõe metros, o que é coerente com v\_P em m/s. O texto também chama o disco de 'cilindro A' em um trecho. Nada disso altera o resultado, porque $\Delta$$\theta$\_A = $\frac12$$\cdot$(a\_t/R\_A)$\cdot$t$^2$ só depende da razão a\_t/R\_A.%
\end{itemize}%
\section*{Questão 23}%
\begin{longtable}{@{}p{0.06\textwidth} p{0.33\textwidth} p{0.33\textwidth} p{0.1\textwidth} p{0.08\textwidth}@{}}\toprule Item & Resolução & Gabarito & Veredito & Certo \\ \midrule
23 & (B) 0,50 V & B & igual & ambos \\
\bottomrule\end{longtable}%
\begin{itemize}%
\item%
\textbf{23)} Confirmei pelo original (pág. 3 do PDF, prova AMARELA, p. 10/18). O máximo do gráfico está em R = 10 $\Omega$ com P = 0,025 W. Como o máximo ocorre em R = R0, temos R0 = 10 $\Omega$. Então U$^2$ = P$\cdot$R = 0,025$\cdot$10 = 0,25 V$^2$, e U = 0,50 V. A conta se confirma por dois caminhos. Primeiro, $\varepsilon$ = 2U = 1,0 V dá Pmax = $\varepsilon$$^2$/(4R0) = 0,025 W. Segundo, em R = 50 $\Omega$, P = 50/3600 $\approx$ 0,0139 W, o que bate com o fim da curva ($\approx$0,014 W). A resolução e o gabarito dão a mesma resposta, e ela está correta.%
\end{itemize}%
\paragraph{Pontos de dúvida}%
\begin{itemize}%
\item%
\textit{(statement)} O valor de Pmax = 0,025 W e o valor de R = 10 $\Omega$ no ponto de máximo foram lidos no gráfico. Na figura, o pico está bem sobre a linha de grade de 0,025 W e sobre R = 10 $\Omega$, e o ponto em 50 $\Omega$ ($\approx$0,014 W) também confere com $\varepsilon$ = 1,0 V. A leitura parece segura.%
\end{itemize}%
\section*{Questão 24}%
\begin{longtable}{@{}p{0.06\textwidth} p{0.33\textwidth} p{0.33\textwidth} p{0.1\textwidth} p{0.08\textwidth}@{}}\toprule Item & Resolução & Gabarito & Veredito & Certo \\ \midrule
Questão 24 (única) & (E) Linha 1; F\_(y) no sentido negativo de y (F\_x = 4Mg sen$\theta$ cos$\theta$ > 0; F\_y = Mg(1 - 4cos$^2$$\theta$) < 0, porque $\theta$ < 60$^\circ$) & E & igual & ambos \\
\bottomrule\end{longtable}%
\begin{itemize}%
\item%
\textbf{Questão 24 (única))} Refiz as contas. Tomando momentos em A, com a normal da estaca perpendicular à barra e aplicada a L/4 de A: N$\cdot$L/4 = Mg$\cdot$L cos$\theta$, logo N = 4Mg cos$\theta$. Da resultante nula: F\_x = N sen$\theta$ > 0 e F\_y = Mg - N cos$\theta$ = Mg(1 - 4cos$^2$$\theta$). F\_y é negativo para $\theta$ < 60$^\circ$. Na figura, $\theta$ fica entre 36$^\circ$ e 42$^\circ$, então F aponta para a direita e para baixo, ao longo da linha 1, que desce da esquerda para a direita passando por A. Conferi numericamente com $\theta$ = 37$^\circ$: N $\approx$ 64 N, F\_x $\approx$ 38,5 N, F\_y $\approx$ -31 N, |F| $\approx$ 49 N, a cerca de 39$^\circ$ abaixo da horizontal. Também conferi a identidade |F|$^2$ = (Mg)$^2$(1 + 8cos$^2$$\theta$) e o ponto de concorrência Q do teorema das três forças: ambos estão corretos. As eliminações também estão certas: a linha 2 exigiria F\_x = 0; a linha 3 daria momento não nulo de T em relação ao apoio P; a (D) exigiria F\_x < 0. A resolução e o gabarito coincidem e estão corretos.%
\end{itemize}%
\paragraph{Pontos de dúvida}%
\begin{itemize}%
\item%
\textit{(statement)} A distância L/4 entre A e a estaca só aparece como cota na figura 1, não no texto. No desenho, o topo da estaca parece mais longe de A do que L/4 (perto de L/2). Isso não muda a resposta: com o apoio a uma distância d, F\_y = Mg(1 - L cos$^2$$\theta$/d). Mesmo com d = L/2 basta $\theta$ < 45$^\circ$, e a figura mostra $\theta$ $\approx$ 36$^\circ$--42$^\circ$.%
\item%
\textit{(solution)} A resolução supõe contato liso no topo da estaca, com a normal perpendicular à barra. Com atrito no apoio, o problema ficaria estaticamente indeterminado. É a hipótese padrão e foi declarada.%
\item%
\textit{(solution)} O ângulo $\theta$ $\approx$ 35$^\circ$--40$^\circ$ foi lido na figura. A conclusão só depende de $\theta$ < 60$^\circ$, uma margem confortável.%
\item%
\textit{(statement)} Os dados M = 2,0 kg e g = 10 m/s$^2$ não são necessários para uma resposta qualitativa. A resolução só os usa para estimar ordens de grandeza.%
\end{itemize}%
\section*{Questão 25}%
\begin{longtable}{@{}p{0.06\textwidth} p{0.33\textwidth} p{0.33\textwidth} p{0.1\textwidth} p{0.08\textwidth}@{}}\toprule Item & Resolução & Gabarito & Veredito & Certo \\ \midrule
25 & (D) O processo A é isobárico, pois tem a maior inclinação nC e Cp > Cv. Nele W = p(V2 - V1) = nR(T2 - T1) > 0, qualquer que seja a atomicidade do gás. & D & igual & ambos \\
\bottomrule\end{longtable}%
\begin{itemize}%
\item%
\textbf{25)} Conferi a transcrição com a página 4 do PDF (prova AMARELA, Questão 25) e ela bate com o original. Refazendo: Q = nC(T - T1), então no gráfico Q$\times$T a inclinação é nC. A reta A, mais inclinada, tem C = Cp e é o processo isobárico; a reta B tem C = Cv e é o isovolumétrico (Mayer: Cp = Cv + R). (A) é falsa: p2 = p1 < p3 = p1$\cdot$T2/T1. (B) é falsa: no processo A o gás realiza trabalho. (C) é falsa: $\Delta$U = nCv$\Delta$T depende da atomicidade. (D) é verdadeira: W = p1(V2 - V1) = nR(T2 - T1) > 0, resultado que só usa pV = nRT. (E) é falsa: U depende só de T, então U2 = U3. Resolução e gabarito marcam (D), e ambos estão certos.%
\end{itemize}%
\paragraph{Pontos de dúvida}%
\begin{itemize}%
\item%
\textit{(solution)} A identificação de qual processo é o isobárico vem só da leitura qualitativa do gráfico: a reta A (contínua) é visivelmente mais inclinada que a B. Também se supõe capacidade calorífica constante, o que está implícito nas retas do gráfico. As duas hipóteses são razoáveis e coerentes com a figura original.%
\item%
\textit{(solution)} No item (E), a resolução escreve U = nCvT. Isso pressupõe Cv constante e U = 0 em T = 0. Para o argumento basta dizer que U de um gás ideal depende só de T, e é isso que a resolução usa de fato. É uma imprecisão de redação, não muda nada.%
\end{itemize}%
\section*{Questão 26}%
\begin{longtable}{@{}p{0.06\textwidth} p{0.33\textwidth} p{0.33\textwidth} p{0.1\textwidth} p{0.08\textwidth}@{}}\toprule Item & Resolução & Gabarito & Veredito & Certo \\ \midrule
26 & (C) 0,5 & C & igual & ambos \\
\bottomrule\end{longtable}%
\begin{itemize}%
\item%
\textbf{26)} A transcrição está de acordo com a p. 11/18 da prova AMARELA. Refiz as contas. Com o gás diatômico, $\gamma$ = 7/5. O calor recebido é Q\_AB = W\_AB = 500nR$\cdot$ln(5/4)/(2/5) = 1250nR$\cdot$ln(5/4) $\approx$ 278,9nR. O calor cedido é |Q\_BC| = (5/2)nR$\cdot$100 = 250nR. Então $\eta$\_M = 1 - 1/(5 ln(5/4)) $\approx$ 1 - 0,90 = 0,10 (sem arredondar, 0,1037). O rendimento de Carnot é $\eta$\_máx = 1 - 400/500 = 0,20. A razão dá 0,5 (sem arredondar, 0,519), e a alternativa C é a única compatível. A expressão da área dada no enunciado é coerente com a adiabática CA: (V\_B/V\_A)\^{}($\gamma$-1) = T\_A/T\_C = 5/4.%
\end{itemize}%
\paragraph{Pontos de dúvida}%
\begin{itemize}%
\item%
\textit{(statement)} O enunciado não diz que o gás é ideal. Também não diz que 'diatômico' significa C\_V = 5R/2 (sem graus de vibração). Isso é implícito, e a expressão dada da área só é coerente com $\gamma$ = 7/5 junto com a adiabática. A resolução declara essas hipóteses.%
\item%
\textit{(statement)} Com o dado 1/(5 ln(5/4)) $\approx$ 0,90, a razão sai exatamente 0,5. Sem arredondar, ela vale $\approx$ 0,519. A diferença não muda a alternativa, porque a mais próxima depois de C é 0,9.%
\end{itemize}%
\section*{Questão 27}%
\begin{longtable}{@{}p{0.06\textwidth} p{0.33\textwidth} p{0.33\textwidth} p{0.1\textwidth} p{0.08\textwidth}@{}}\toprule Item & Resolução & Gabarito & Veredito & Certo \\ \midrule
Q27 & (A) y(x) = 3x - 9/2 & A & igual & ambos \\
\bottomrule\end{longtable}%
\begin{itemize}%
\item%
\textbf{Q27)} Conferi o enunciado na página original (pág. 12/18 da prova AMARELA) e ele bate com a transcrição. O plano xy é horizontal e não tem atrito. Peso e normal são verticais e se anulam, e a força da mola é interna, então a resultante externa no plano é nula e v\_CM é constante. Como as massas são iguais, r\_CM(0) = ((0,3)+(4,0))/2 = (2; 1,5) m e v\_CM = (1î+3?)/2 = (0,5; 1,5) m/s. Daí y - 1,5 = (1,5/0,5)(x - 2), ou seja, y = 3x - 4,5 = 3x - 9/2. A resolução e o gabarito dão a alternativa (A), e as duas estão certas. O valor de k não interfere no resultado.%
\end{itemize}%
\paragraph{Pontos de dúvida}%
\begin{itemize}%
\item%
\textit{(solution)} Erro de digitação no LaTeX: '\textbackslash{}emph\{interna*\}' tem um asterisco sobrando. É só cosmético.%
\item%
\textit{(statement)} A reta y = 3x - 9/2 só vale para x $\geq$ 2 m (t $\geq$ 0), porque fora disso a trajetória vira uma semirreta. A resolução diz isso, mas a alternativa (A) não traz a restrição. Isso não muda a resposta.%
\end{itemize}%
\section*{Questão 28}%
\begin{longtable}{@{}p{0.06\textwidth} p{0.33\textwidth} p{0.33\textwidth} p{0.1\textwidth} p{0.08\textwidth}@{}}\toprule Item & Resolução & Gabarito & Veredito & Certo \\ \midrule
28 & (B) E = -($\pi$$^2$ M k$^2$ q$^{4}$ / T$^2$)\^{}\{1/3\} & B & igual & ambos \\
\bottomrule\end{longtable}%
\begin{itemize}%
\item%
\textbf{28)} Enunciado conferido com o PDF original (p. 12/18, prova AMARELA), sem diferenças. Refiz a conta: como as massas são iguais, o CM fica no ponto médio e cada partícula gira com raio d/2. M$\omega$$^2$(d/2) = 2kq$^2$/d$^2$ $\Rightarrow$ d$^3$ = kq$^2$T$^2$/($\pi$$^2$M). Ec = M$\omega$$^2$d$^2$/4 = kq$^2$/d e Ep = -2kq$^2$/d, logo E = -kq$^2$/d (virial). Então E$^3$ = -k$^3$q$^{6}$$\cdot$$\pi$$^2$M/(kq$^2$T$^2$) = -$\pi$$^2$Mk$^2$q$^{4}$/T$^2$, ou seja, E = -($\pi$$^2$Mk$^2$q$^{4}$/T$^2$)\^{}\{1/3\}. Conferi numericamente (k = 8,99e9, q = 1,3 µC, M = 0,7 kg, T = 2,3 s): a energia calculada direto, -0,067052 J, bate com a fórmula (B). A (D) dá -0,1064 J. A análise dimensional que descarta (A), (C) e (E) está correta. A resolução e o gabarito estão certos.%
\end{itemize}%
\paragraph{Pontos de dúvida}%
\begin{itemize}%
\item%
\textit{(solution)} A resolução diz que a alternativa (D) sai de usar o raio d no lugar de d/2. A conta não confirma isso. Com M$\omega$$^2$d = 2kq$^2$/d$^2$ e E = -kq$^2$/d, o resultado é E = -(2$\pi$$^2$Mk$^2$q$^{4}$/T$^2$)\^{}\{1/3\}, com fator 2 e não 4. Se, além disso, a energia cinética for calculada com raio d, dá E = 0. Para chegar ao fator 4 de (D) é preciso, por exemplo, usar d$^3$ = kq$^2$/(M$\omega$$^2$), o que exige combinar dois erros (raio d e carga q em vez de 2q). É só um comentário sobre o distrator, sem efeito na resposta final.%
\end{itemize}%
\section*{Questão 29}%
\begin{longtable}{@{}p{0.06\textwidth} p{0.33\textwidth} p{0.33\textwidth} p{0.1\textwidth} p{0.08\textwidth}@{}}\toprule Item & Resolução & Gabarito & Veredito & Certo \\ \midrule
Q29 & (A) F = (1/10)$\cdot$GM$^2$/R$^2$ & A & igual & ambos \\
\bottomrule\end{longtable}%
\begin{itemize}%
\item%
\textbf{Q29)} O enunciado e a figura da página 12/18 da prova (página 5 do PDF, prova AMARELA) conferem com a transcrição. A cavidade tem raio R/2 e centro a R/2 à esquerda de O. A partícula está a 2R à direita de O, então fica a 5R/2 do centro da cavidade. A massa retirada é M/8. Pela superposição, F = GMm/R$^2$$\cdot$(1/4 - (1/8)(4/25)) = (23/100)$\cdot$GMm/R$^2$. Com m = 10M/23, F = GM$^2$/(10R$^2$). Conferi a conta com python (frações exatas): 23/100 e 1/10. A resolução e o gabarito estão certos.%
\end{itemize}%
\paragraph{Pontos de dúvida}%
\begin{itemize}%
\item%
\textit{(solution)} A resolução usa o teorema das cascas para a esfera de massa negativa (cavidade), o que exige a partícula fora da cavidade. Isso vale (5R/2 > R/2) e a resolução declara a hipótese. A conta de ordem de grandeza (10/92 $\approx$ 0,109) também está correta. Não há ressalvas relevantes.%
\end{itemize}%
\section*{Questão 30}%
\begin{longtable}{@{}p{0.06\textwidth} p{0.33\textwidth} p{0.33\textwidth} p{0.1\textwidth} p{0.08\textwidth}@{}}\toprule Item & Resolução & Gabarito & Veredito & Certo \\ \midrule
30 & (D) 6,0$\times$10\^{}{-}7 m & D & igual & ambos \\
\bottomrule\end{longtable}%
\begin{itemize}%
\item%
\textbf{30)} Conferi com o original (pág. 13/18, prova AMARELA). No gráfico, y3 está em 10\^{}{-}4 m e L em $\times$1000$\cdot$d m. A variação indicada é $\Delta$L = 0,5$\times$1000d = 500d para $\Delta$y3 = 9,0$\times$10\^{}{-}4 m. Com y3 = 3$\lambda$L/d, fica 9,0$\times$10\^{}{-}4 = 3$\lambda$$\cdot$500, e daí $\lambda$ = 6,0$\times$10\^{}{-}7 m (conferido em Python). Usar variações em vez de valores absolutos é correto. O segmento do gráfico, se prolongado, passa aproximadamente pela origem, o que é coerente com y3 $\propto$ L. A resolução e o gabarito coincidem e estão corretos.%
\end{itemize}%
\paragraph{Pontos de dúvida}%
\begin{itemize}%
\item%
\textit{(solution)} A resolução diz que a alternativa (C) 3,0$\times$10\^{}{-}7 viria de 'contar errado o número de franjas'. Esse valor também sai de ler mal a escala do eixo L ($\Delta$L = 1000d em vez de 500d). Isso é só especulação sobre os distratores e não afeta a resposta.%
\item%
\textit{(statement)} Os valores 0,5 e 9,0 são variações marcadas no próprio gráfico, que não traz valores absolutos nos eixos. A resolução trata isso corretamente usando a inclinação. O enunciado também não diz explicitamente que d fica fixo enquanto L varia, mas a escala 'em múltiplos de d' faz d se cancelar de qualquer forma.%
\end{itemize}%
\section*{Questão 31}%
\begin{longtable}{@{}p{0.06\textwidth} p{0.33\textwidth} p{0.33\textwidth} p{0.1\textwidth} p{0.08\textwidth}@{}}\toprule Item & Resolução & Gabarito & Veredito & Certo \\ \midrule
31 & (C) 0,20 m/s & C & igual & ambos \\
\bottomrule\end{longtable}%
\begin{itemize}%
\item%
\textbf{31)} Refiz a conta. Com B(t) = 2,0t (em módulo) e A(t) = 0,10$\cdot$(0,20 + vt), o fluxo é ? = 0,040t + 0,20vt$^2$, e então $\varepsilon$ = d?/dt = 0,040 + 0,40vt. Em t = 2,0 s, $\varepsilon$ = RI = 4,0$\cdot$0,050 = 0,20 V, o que dá 0,040 + 0,80v = 0,20 e v = 0,20 m/s. Conferi também com derivada numérica em Python, que dá I = 0,050 A. As parcelas A$\cdot$dB/dt = 0,12 V e BLv = 0,080 V têm o mesmo sinal e se somam, como a resolução diz. Pela lei de Lenz o sentido sai anti{-}horário visto de +z, o que também está correto. As alternativas{-}armadilha que a resolução cita (0,50 e 0,40 m/s) também conferem. Comparei a transcrição com a página 13/18 do PDF: o enunciado, os dados (L = 0,10 m, R = 4,0 $\Omega$, x(0) = 0,20 m, B = -(2,0t)k?, I = 0,050 A em t = 2,0 s) e as alternativas batem.%
\end{itemize}%
\paragraph{Pontos de dúvida}%
\begin{itemize}%
\item%
\textit{(statement)} O enunciado diz que o comprimento da barra é L = 0,10 m, e a resolução usa esse valor como distância entre os trilhos. A figura mostra trilhos em y = 0 e y = 0,1 m, então a hipótese é coerente, mas fica implícita no texto.%
\end{itemize}%
\section*{Questão 32}%
\begin{longtable}{@{}p{0.06\textwidth} p{0.33\textwidth} p{0.33\textwidth} p{0.1\textwidth} p{0.08\textwidth}@{}}\toprule Item & Resolução & Gabarito & Veredito & Certo \\ \midrule
32 & (E) T/$\mu$ = (0,60/t')$^2$ & E & igual & ambos \\
\bottomrule\end{longtable}%
\begin{itemize}%
\item%
\textbf{32)} Pelo original (pág. 14/18 da prova AMARELA), A (grossa) tem máximo em 0, mínimo perto de t' e máximo perto de 2t', logo T\_A = 2t'. B (pontilhada) faz 3 ciclos e C (tracejada) faz 5 ciclos em 2t'. R repete com período 2t'. Como f\_R = f0, o fundamental é A, com f0 = 1/(2t'), f\_B = 3f0 e f\_C = 5f0. Com corda fixa nas duas pontas, $\lambda$0 = 2L = 1,20 m e v = $\lambda$0$\cdot$f0 = 0,60/t', de modo que T/$\mu$ = v$^2$ = (0,60/t')$^2$. Isso torna (E) verdadeira. As outras são falsas: (A) os harmônicos são 1, 3 e 5, não 1, 4 e 6; (B) o fundamental é A, não C; (C) f0 = 1/(2t'), não 1/t'; (D) f\_B/f\_C = 3/5, não 2/5. Só (E) é correta. Conferi as contas em python.%
\end{itemize}%
\paragraph{Pontos de dúvida}%
\begin{itemize}%
\item%
\textit{(statement)} Os períodos saem da contagem de ciclos no gráfico, que é esquemático. Na figura, o máximo da curva A perto de 2t' parece um pouco à direita da marca 2t', onde está o pico de R. É só imprecisão do desenho. Não muda a identificação dos harmônicos 1, 3 e 5, que também se confirma pelo mínimo profundo de R em t' (todas as componentes ímpares em mínimo).%
\item%
\textit{(solution)} A resolução supõe sem dizer que as componentes têm a forma cos(2$\pi$n f0 t), ou seja, fase nula com máximo em t = 0. A figura sustenta essa leitura, e a hipótese só entra na verificação de coerência, não na resposta.%
\end{itemize}%
\section*{Questão 33}%
\begin{longtable}{@{}p{0.06\textwidth} p{0.33\textwidth} p{0.33\textwidth} p{0.1\textwidth} p{0.08\textwidth}@{}}\toprule Item & Resolução & Gabarito & Veredito & Certo \\ \midrule
33 & (C): com f1 = f2 + 40 Hz há batimentos em qualquer posição de O. O som percebido tem frequência (f1+f2)/2 = f2 + 20 Hz e a frequência de batimento é |f1-f2| = 40 Hz. & C & igual & ambos \\
\bottomrule\end{longtable}%
\begin{itemize}%
\item%
\textbf{33)} Conferi o enunciado na página 7 (Questão 33) e a transcrição está fiel. Refiz a conta: em O, y1 + y2 = 2A$\cdot$cos($\pi$$\Delta$f$\cdot$t - $\varphi$?)$\cdot$sin($\pi$(f1+f2)t - $\varphi$?). A oscilação rápida tem frequência (f1+f2)/2 = f2 + 20 Hz. A intensidade varia com cos$^2$, então os batimentos têm frequência $\Delta$f = 40 Hz. A posição de O só aparece na fase constante $\varphi$?, por isso há batimentos em todas as posições e a (C) é verdadeira. A (D) está errada porque f2 + 20 Hz é a frequência do som, não a dos batimentos, e eles não ficam restritos a algumas posições. A (E) está errada porque a frequência de batimento é 40 Hz, não 20 Hz, e também vale para todas as posições. Com f1 = f2 forma{-}se uma onda estacionária, com amplitude 2A|cos($\pi$(d2-d1)/$\lambda$)|. Os nós existem porque $\lambda$/2 < d1 + d2, mas fora deles há som, então a (A) é falsa. A frequência ouvida é f1, não 2f1, então a (B) é falsa. A resolução e o gabarito estão certos.%
\end{itemize}%
\paragraph{Pontos de dúvida}%
\begin{itemize}%
\item%
\textit{(statement)} Um batimento de 40 Hz fica acima do limite em que o ouvido distingue batimentos com clareza (cerca de 10--20 Hz) e na prática soa como 'aspereza'. O enunciado segue o modelo de livro{-}texto, então isso não muda a resposta (C), mas pode ser comentado com os alunos.%
\item%
\textit{(solution)} A resolução supõe que as fontes emitem em fase em t0 = 0 (fase inicial nula). Isso só afeta a posição dos nós no caso f1 = f2. A conclusão sobre as alternativas não depende dessa hipótese.%
\item%
\textit{(key)} O gabarito é preliminar (SSPM, 01/09/2026). A resposta C é fisicamente correta, mas é bom conferir se o gabarito definitivo manteve essa alternativa.%
\end{itemize}%
\section*{Questão 34}%
\begin{longtable}{@{}p{0.06\textwidth} p{0.33\textwidth} p{0.33\textwidth} p{0.1\textwidth} p{0.08\textwidth}@{}}\toprule Item & Resolução & Gabarito & Veredito & Certo \\ \midrule
34 & (B) (F)(F)(F): k = 20 N/m, d\_máx = 30 cm, T = 2$\pi$$\surd$(3/20) $\approx$ 2,43 s & (E) (F)(V)(V) & diferente & resolução \\
\bottomrule\end{longtable}%
\begin{itemize}%
\item%
\textbf{34)} Reli a figura 2 em alta resolução. A parábola passa exatamente por (0,05 m; 0,025 J) e (0,10 m; 0,10 J) e termina em U = 0,2 J com d $\approx$ 0,14 m. Logo U = 10d$^2$ e k = 20 N/m, sem ambiguidade. A mola está presa em m1 e o piso não tem atrito, então a única força horizontal sobre m2 é o atrito de m1. Assim, a\_máx = $\mu$e$\cdot$m1$\cdot$g/m2 = 0,4$\cdot$10/2 = 2 m/s$^2$, e k$\cdot$d\_máx = (m1+m2)$\cdot$a\_máx = 6 N, o que dá d\_máx = 0,30 m = 30 cm. As afirmativas I (80 cm) e II (20 cm) são, portanto, falsas. O período do OHS do conjunto é T = 2$\pi$$\surd$((m1+m2)/k) = 2$\pi$$\surd$0,15 $\approx$ 2,43 s. Já 0,4$\pi$$\surd$5 = 2$\pi$$\surd$0,2 $\approx$ 2,81 s, o que exigiria M/k = 0,2, ou seja, k = 15 N/m com M = 3 kg (incompatível com o gráfico) ou M = 4 kg com k = 20. Logo a afirmativa III também é falsa. Conferi as contas em python. O gabarito E só se explica com dois erros juntos: (a) a conta ingênua k$\cdot$d = $\mu$e$\cdot$m1$\cdot$g (d = 4/20 = 0,20 m), que supõe que o atrito equilibra toda a força da mola, quando ele só precisa acelerar m2; e (b) um período calculado com M/k = 0,2, que não sai dos dados impressos. Com m1 = 1,0 kg e m2 = 2,0 kg, a resposta correta é B.%
\end{itemize}%
\paragraph{Pontos de dúvida}%
\begin{itemize}%
\item%
\textit{(key)} O gabarito preliminar (E) marca como verdadeiros d\_máx = 20 cm e T = 0,4$\pi$$\surd$5. Com os dados impressos (k = 20 N/m pelo gráfico, m1 = 1 kg, m2 = 2 kg, $\mu$e = 0,4), d\_máx = 30 cm e T = 2$\pi$$\surd$0,15 $\approx$ 2,43 s $\neq$ 2,81 s. O valor de 20 cm sai da análise errada k$\cdot$d = $\mu$e$\cdot$m1$\cdot$g, que ignora que o atrito só precisa acelerar m2. O período 0,4$\pi$$\surd$5 exigiria (m1+m2)/k = 0,2. Tudo indica erro no gabarito preliminar, e vale conferir se o definitivo alterou ou anulou a questão.%
\item%
\textit{(statement)} Possível erro de digitação no enunciado: com m1 = m2 = 2,0 kg (M = 4 kg) e k = 20 N/m, a física correta dá d\_máx = 0,4$\cdot$20$\cdot$4/(20$\cdot$2) = 0,80 m = 80 cm e T = 2$\pi$$\surd$(4/20) = 0,4$\pi$$\surd$5, ou seja, a alternativa A (V)(F)(V), exatamente. Isso sugere que a questão pode ter sido elaborada com outros valores de massa. Mesmo assim, nenhuma combinação natural de massas leva à E com a análise correta. Com os dados impressos, a resposta é B.%
\item%
\textit{(statement)} O valor de k depende de leitura do gráfico. Os pontos (0,10; 0,10) e (0,05; 0,025) caem exatamente na grade, então k = 20 N/m é seguro. A afirmativa III não traz unidade (0,4$\pi$$\surd$5, presumivelmente em segundos).%
\end{itemize}%
\section*{Questão 35}%
\begin{longtable}{@{}p{0.06\textwidth} p{0.33\textwidth} p{0.33\textwidth} p{0.1\textwidth} p{0.08\textwidth}@{}}\toprule Item & Resolução & Gabarito & Veredito & Certo \\ \midrule
35 & (D) 2,0 m & D & igual & ambos \\
\bottomrule\end{longtable}%
\begin{itemize}%
\item%
\textbf{35)} Conferi a transcrição com o original (página 15/18 da prova AMARELA, questão 35): os dados, a figura e as alternativas batem. Pelo Bernoulli entre a superfície livre e C, a velocidade no tubo é v$^2$ = 2g(H-h). Entre um ponto do tubo e C, com a mesma v, a pressão interna é p(y) = p0 - $\rho$g(y-h). Ela é mínima em B, onde y = H+a, e lá p0 - p\_B = $\rho$g(H+a-h). Impondo $\rho$g(H+a-h) $\leq$ 25 kPa, sai H+a-h $\leq$ 2,5 m, ou seja, h $\geq$ 3,0+1,5-2,5 = 2,0 m. Refiz as contas em Python. As diferenças de pressão para h = 0,9; 1,5; 1,9; 2,0 e 2,5 m dão 36, 30, 26, 25 e 20 kPa, o que confirma a análise das alternativas feita na resolução. A verificação do trecho submerso também está certa: lá a diferença é constante e vale $\rho$g(H-h), menor que em B. A resolução e o gabarito concordam na alternativa D.%
\end{itemize}%
\paragraph{Pontos de dúvida}%
\begin{itemize}%
\item%
\textit{(solution)} A resolução supõe que o tubo tem seção uniforme, para que a velocidade em B seja a mesma da saída C. O enunciado diz apenas 'sifão ideal'. A hipótese é a usual e está declarada, mas, se a seção variasse, a pressão em B mudaria.%
\item%
\textit{(statement)} O tubo colapsa se a diferença 'exceder' 25,0 kPa. Em h = 2,0 m a diferença é exatamente 25 kPa, que é o caso{-}limite e ainda não colapsa. Por isso o mínimo é 2,0 m, o que está de acordo com o gabarito.%
\end{itemize}%
\section*{Questão 36}%
\begin{longtable}{@{}p{0.06\textwidth} p{0.33\textwidth} p{0.33\textwidth} p{0.1\textwidth} p{0.08\textwidth}@{}}\toprule Item & Resolução & Gabarito & Veredito & Certo \\ \midrule
36 & (B) V1 > V2 quando d1 > d2 (V1/V2 = d1/d2, V1 + V2 = V) & B & igual & ambos \\
\bottomrule\end{longtable}%
\begin{itemize}%
\item%
\textbf{36)} A transcrição confere com o original (p. 16/18 da prova AMARELA, questão 36). Refiz a física: a placa metálica neutra e isolada só se polariza (-$\sigma$ em cima, +$\sigma$ embaixo). Por isso o campo vale E = $\sigma$/$\varepsilon$0 nas duas regiões de ar, o que equivale a dois capacitores em série com a mesma carga Q. Como E é o mesmo nas duas regiões, V1 = E$\cdot$d1, V2 = E$\cdot$d2 e V1 + V2 = V, valor fixado pela bateria e independente da posição da placa. Daí E = V/(d1+d2) e V1/V2 = d1/d2: se d1 > d2, então V1 > V2, que é a alternativa (B). A alternativa (A) é falsa, salvo quando d1 = d2. (C) é o oposto do resultado. (D) é falsa porque a soma é V. (E) é falsa porque o campo só é nulo dentro do metal. A resolução e o gabarito estão certos.%
\end{itemize}%
\paragraph{Pontos de dúvida}%
\begin{itemize}%
\item%
\textit{(solution)} Na análise de (D), a resolução diz que só a divisão entre V1 e V2 muda 'porque d1+d2 é constante'. Isso vale porque a espessura da placa metálica é fixa, hipótese implícita mas razoável. Não altera a conclusão.%
\end{itemize}%
\section*{Questão 37}%
\begin{longtable}{@{}p{0.06\textwidth} p{0.33\textwidth} p{0.33\textwidth} p{0.1\textwidth} p{0.08\textwidth}@{}}\toprule Item & Resolução & Gabarito & Veredito & Certo \\ \midrule
37 & (B) 6,0$\cdot$10\^{}{-}1 J (W\_Nx = mgh$\cdot$cos$^2$30$^\circ$ = 0,60 J) & B & igual & ambos \\
\bottomrule\end{longtable}%
\begin{itemize}%
\item%
\textbf{37)} Conferi com a figura original (pág. 16/18 da prova AMARELA). O plano desce para a direita, com o ângulo de 30$^\circ$ no canto inferior direito, e x aponta para a direita. A normal aponta para cima e para a direita: N\_x = mg$\cdot$cos30$^\circ$$\cdot$sen30$^\circ$ > 0. O deslocamento é d = h/sen30$^\circ$ = 0,40 m, com $\Delta$x = d$\cdot$cos30$^\circ$ > 0. Assim, W = mg$\cdot$cos$^2$30$^\circ$$\cdot$h = 0,4$\cdot$10$\cdot$0,2$\cdot$0,75 = 0,60 J. Conferi em python: W\_Nx = +0,60 J e W\_Ny = -0,60 J, o que dá trabalho total nulo, como esperado. A resolução e o gabarito coincidem na alternativa B.%
\end{itemize}%
\paragraph{Pontos de dúvida}%
\begin{itemize}%
\item%
\textit{(statement)} A conta supõe que o bloco desliza apoiado no plano e que o plano está fixo, como diz o enunciado. Com isso a normal é constante e vale mg$\cdot$cos30$^\circ$. O sinal positivo de W\_Nx depende de ler na figura que o plano desce para a direita e que x aponta para a direita. A figura original confirma as duas coisas.%
\end{itemize}%
\section*{Questão 38}%
\begin{longtable}{@{}p{0.06\textwidth} p{0.33\textwidth} p{0.33\textwidth} p{0.1\textwidth} p{0.08\textwidth}@{}}\toprule Item & Resolução & Gabarito & Veredito & Certo \\ \midrule
38 & (C) x = (15/8)$\cdot$mg/k & C & igual & ambos \\
\bottomrule\end{longtable}%
\begin{itemize}%
\item%
\textbf{38)} Conferi o enunciado na p. 17/18 do PDF (página 10 do arquivo). Ele confere com a transcrição: L/3 à esquerda de O, 2L/3 à direita, $\rho$c = 8$\rho$ e metade do volume submerso. Refiz a conta. O volume é V = m/$\rho$c = m/(8$\rho$), então o empuxo é E = $\rho$$\cdot$g$\cdot$V/2 = mg/16. A tração no fio é T = mg - mg/16 = 15mg/16. Tomando os torques em torno de O: kx$\cdot$(L/3) = T$\cdot$(2L/3), logo kx = 2T = 15mg/8 e x = (15/8)$\cdot$mg/k, que é a alternativa (C). A resolução e o gabarito batem e estão certos. Os casos{-}limite da resolução também conferem: sem empuxo, x = 2mg/k; com o cilindro todo submerso, x = 7mg/(4k).%
\end{itemize}%
\paragraph{Pontos de dúvida}%
\begin{itemize}%
\item%
\textit{(solution)} A resolução supõe que o cilindro não toca o fundo do recipiente e que o peso do fio é desprezível. A figura confirma isso, porque o cilindro aparece suspenso. A resolução declara a hipótese, e ela não muda o resultado.%
\end{itemize}%
\section*{Questão 39}%
\begin{longtable}{@{}p{0.06\textwidth} p{0.33\textwidth} p{0.33\textwidth} p{0.1\textwidth} p{0.08\textwidth}@{}}\toprule Item & Resolução & Gabarito & Veredito & Certo \\ \midrule
39 & (E) C\_B$\cdot$(T\_eq - T\_i,B)/{[}C\_A$\cdot$|T\_eq - T\_i,A|{]} & E & igual & ambos \\
\bottomrule\end{longtable}%
\begin{itemize}%
\item%
\textbf{39)} Conferi o enunciado e a figura no PDF original (p. 17/18 da prova AMARELA) e a transcrição está fiel. Refiz a conta. Com taxa P constante durante todo o intervalo {[}0, t2{]}, A cede C\_A(T\_i,A - T\_eq) = P$\cdot$t2. De 0 a t1 a temperatura de B fica constante: é o patamar de fusão, e esse calor é latente. De t1 a t2, B aquece com C\_B(T\_eq - T\_i,B) = P(t2 - t1). Dividindo uma equação pela outra, P se cancela e (t2 - t1)/t2 = C\_B(T\_eq - T\_i,B)/{[}C\_A(T\_i,A - T\_eq){]}. Como T\_i,A > T\_eq, isso é igual à alternativa (E). A resolução descarta as alternativas A, C e D porque usam T\_eq isolada, não uma diferença de temperaturas, e o resultado dependeria da escala. Descarta B porque troca as variações de temperatura entre os corpos. Os dois argumentos estão corretos.%
\end{itemize}%
\paragraph{Pontos de dúvida}%
\begin{itemize}%
\item%
\textit{(statement)} O enunciado chama C\_B de 'capacidade térmica da substância B', que entra no recipiente sólida. Mas a conta só usa o aquecimento depois da fusão, entre t1 e t2, quando B já está líquida. Então C\_B tem que ser a capacidade térmica de B na fase líquida. A resolução avisa isso ('C\_B é a capacidade térmica de B nessa etapa'), mas o enunciado não deixa claro. Isso não muda a alternativa.%
\item%
\textit{(solution)} A resolução diz que o patamar de B entre 0 e t1 é uma fusão. Essa leitura é a única coerente com B absorver calor a temperatura constante. Também supõe que A não muda de fase, o que é coerente com a reta contínua sem patamar. As duas hipóteses saem da figura e são razoáveis.%
\end{itemize}%
\section*{Questão 40}%
\begin{longtable}{@{}p{0.06\textwidth} p{0.33\textwidth} p{0.33\textwidth} p{0.1\textwidth} p{0.08\textwidth}@{}}\toprule Item & Resolução & Gabarito & Veredito & Certo \\ \midrule
40 & (D) v\_vb/v\_vm = $\surd$6/2 $\approx$ 1,22 (com v\_vm = 10($\surd$3+1) $\approx$ 27,3 km/h e v\_vb $\approx$ 33,5 km/h) & D & igual & ambos \\
\bottomrule\end{longtable}%
\begin{itemize}%
\item%
\textbf{40)} Refiz a conta com v\_vb = v\_vm - v\_bm, tomando x para leste e y para norte, com v\_bm = 10 x?. O barco não tem velocidade na direção norte, então a componente norte do vento é a mesma nos dois referenciais: v\_vb$\cdot$cos45$^\circ$ = v\_vm$\cdot$cos30$^\circ$. Daí a razão é $\surd$6/2. A equação da componente leste fica v\_vm/2 + 10 = v\_vm$\cdot$$\surd$3/2, o que dá v\_vm = 10($\surd$3+1) $\approx$ 27,32 km/h. Conferi em Python: o vento relativo ao barco sai com 33,46 km/h, a exatamente 45$^\circ$ do norte para o oeste, e a razão é 1,2247 = $\surd$6/2. A figura do original (pág. 11, 'Página 18/18', prova AMARELA) bate com a transcrição: o traço da bandeira sai do centro do barco para cima e para a esquerda, a 45$^\circ$ do N. A resolução está correta e o gabarito D também.%
\end{itemize}%
\paragraph{Pontos de dúvida}%
\begin{itemize}%
\item%
\textit{(statement)} O enunciado não diz se '30$^\circ$ para o oeste' é a direção para onde o vento sopra ou de onde ele vem. A resolução adota 'para onde sopra' e mostra que a leitura oposta (vento vindo de N30$^\circ$O) levaria a v\_vm < 0 na equação da componente leste, então essa leitura não tem solução física. A razão $\surd$6/2 sai das componentes norte de qualquer forma. Por isso a ambiguidade não muda a alternativa, mas vale citar.%
\item%
\textit{(solution)} A resolução supõe sem discutir que a bandeira aponta no sentido para onde vai o vento aparente (a favor do vento). É a hipótese física padrão e combina com a figura, mas não é dita explicitamente no enunciado.%
\end{itemize}%
\end{document}